\documentclass[11pt]{article}
\usepackage{preamble}
\renewcommand{\answer}[1]{}

\begin{document}

\ladderheading{AP Statistics \textperiodcentered\ Unit 2 \textperiodcentered\ Topic 2.7 \textperiodcentered\ B: 2026-10-29 \textperiodcentered\ E: 2026-11-18}{Overlap Does Not Decide Independence}

\focusbanner{TODAY'S PROBLEM}

Suppose a school surveys 200 students about whether they own at least one pet and whether they play a musical instrument. Let $P$ be the event that a randomly selected surveyed student owns a pet and $I$ the event that the student plays an instrument.\par\smallskip \textbf{Aaliyah:} \emph{The events overlap because 84 students do both. Since they are not mutually exclusive, they must be independent.}\par \textbf{Marcus:} \emph{Overlap does not decide independence; the relevant probabilities must be compared.}\par
\begin{center}
\textbf{Pet ownership and instrument playing}\par\smallskip
\begin{tabular}{|l|c|c|c|}
\hline
\textbf{Pet?} & \textbf{Plays} & \textbf{Does not} & \textbf{Total} \\ \hline
\textbf{Yes} & 84 & 36 & 120 \\ \hline
\textbf{No} & 16 & 64 & 80 \\ \hline
\textbf{Total} & 100 & 100 & 200 \\ \hline
\end{tabular}
\end{center}
\begin{firsttakebox}{FIRST TAKE --- before the video (5 min)}
Does owning a pet seem to change the chance that a student plays an instrument? Give one table detail in one sentence.
\workspace{0.9in}
\end{firsttakebox}

\begin{callout}{\IconBook\ INDEPENDENCE, INTERSECTION, AND UNION (keep on the page)}{calloutgreen}
Events $A$ and $B$ are \textbf{independent} exactly when knowing one does not change the probability of the other:
$P(A\mid B)=P(A)$. Equivalently, $P(A\cap B)=P(A)P(B)$. Establish independence before using that product.
\textbf{Mutually exclusive} events have no overlap, so $P(A\cap B)=0$; this is a different relationship.
For any two events, $P(A\cup B)=P(A)+P(B)-P(A\cap B)$. To calculate a conditional probability from a table, divide the count in both events by the total for the given group.
\end{callout}

\newpage
\textbf{Finish after the video:}\par\smallskip

\dokbadge{1}~\textbf{(a)} State the counts for $P$, $I$, and $P\cap I$.\par
\workspace{0.7in}

\dokbadge{2}~\textbf{(b)} Find $P(I)$, $P(I\mid P)$, and $P(P\cup I)$. Show the fractions or setup.\par
\workspace{1.4in}

\dokbadge{3}~\textbf{(c)} Whose reasoning is supported? Decide independence using the two instrument-playing probabilities. Explain why the 84 students in both groups settle mutual exclusivity but do not settle independence. Write 3--4 sentences.\par
\workspace{2.0in}

\begin{sentenceframebox}\raggedright\textbf{Frame:} The events are \blank[0.9in] because $P(I\mid P)=\blank[0.5in]$ while $P(I)=\blank[0.5in]$.\end{sentenceframebox}

\begin{sentenceframebox}\raggedright\textbf{Frame:} Mutually exclusive means \blank[1.8in], whereas independent means \blank[2.0in].\end{sentenceframebox}

\begin{summaryexitbox}{EXIT --- turn this in}
One thing the video changed about my first take: \hrulefill\par
\workspace{0.4in}
\end{summaryexitbox}

\end{document}
